\chapter{Worked Example: OAuth Authorization Code Flow}
\label{ch:oauth-example}

\index{OAuth}
This chapter is a behavioural sketch inspired by the OAuth authorization-code
flow. It is not a replacement for the OAuth specification and deliberately begins
above HTTP and protobuf. Its purpose is to exercise scenario decomposition,
parameter flow, visible evidence, alternatives, requirements, and rendering.

\section{Actors, values, and observable events}

The example has four participants: a person using a browser, an application, an
authorization service, and a resource service. The browser is modelled separately
from the person because redirects are observable messages between software
participants.

\begin{lstlisting}
inductive Actor where
  | user | browser | app | authorizationServer | resourceServer

opaque State       : Type
opaque AuthCode    : Type
opaque AccessToken : Type
opaque Resource    : Type
opaque Profile     : Type
opaque Password    : Type
opaque Bytes       : Type
\end{lstlisting}

Every terminal elaborates to the common event envelope. For example, the vague
message ``authorization redirected'' becomes conceptually:

\begin{lstlisting}
{ tstamp := t2, id := e2, priorIds := [e1]
, from := Actor.app, to := Actor.browser
, kind := "authorization redirected"
, fields := [field "state" state, field "redirectUri" callback]
}
\end{lstlisting}

\section{Scenario 1: begin authorization}

Before the redirect returns, the authorization service may authenticate the person.
OAuth does not prescribe the authentication mechanism; we include a password branch
to exercise secret input without pretending that the password is an OAuth protocol
parameter.

\begin{lstlisting}
def authenticateAtAuthorizationServer (password : Password) : Scenario :=
  scenario "authenticate resource owner" do
    from_ .browser
    to_ .authorizationServer
    emitBytes "credentials submitted"
      (protectFor .authorizationServer
        (encode [secretField "password" password]))

    from_ .authorizationServer
    to_ .browser
    oneOf "authentication result"
      [ emitBytes "authentication accepted" (encode [tag 1])
      , emitBytes "authentication rejected" (encode [tag 0])
      ]
\end{lstlisting}

The wire bytes are specified even though the transport is not. The abstract
\texttt{protectFor} constructor promises confidentiality to the named endpoint. A
later refinement may realize that promise with TLS, message-level encryption, or a
local trusted channel. Merely placing the password inside a field would not support
a secrecy theorem: a network observer would then derive it from the bytes.

\begin{lstlisting}
def beginAuthorization (resource : Resource) : Scenario :=
  scenario "begin authorization" do
    let state <- declare State "state"
    let callback := Uri.https "app.example" "/oauth/callback"

    from_ .user
    to_ .browser
    emit "resource requested" [field "resource" resource]

    from_ .browser
    to_ .app
    emit "application request" [field "resource" resource]

    from_ .app
    to_ .browser
    emit "authorization redirected"
      [field "state" state, field "redirectUri" callback]

    from_ .browser
    to_ .authorizationServer
    emit "authorization requested"
      [field "state" state, field "redirectUri" callback]

    export "state" state
    export "resource" resource
\end{lstlisting}

This scenario establishes the correlation value \texttt{state}. Nothing here says
how the application stores it. The contract only requires later observable
messages to carry the corresponding value.

\section{Scenario 2: grant or deny}

\begin{lstlisting}
def authorizationGranted : Scenario :=
  scenario "authorization granted" do
    use state <- imported State "state"
    let code <- declare AuthCode "code"

    from_ .authorizationServer
    to_ .browser
    emit "authorization code returned"
      [field "state" state, field "code" code]

    from_ .browser
    to_ .app
    emit "authorization callback"
      [field "state" state, field "code" code]
    export "code" code

def authorizationDenied : Scenario :=
  scenario "authorization denied" do
    use state <- imported State "state"
    from_ .authorizationServer
    to_ .browser
    emit "authorization error returned"
      [field "state" state, field "reason" "access_denied"]
    from_ .browser
    to_ .app
    emit "authorization error callback"
      [field "state" state, field "reason" "access_denied"]

def authorizationDecision : Scenario :=
  oneOf "authorization decision"
    [authorizationGranted, authorizationDenied]
\end{lstlisting}

The abstract contract makes both decisions possible. A product model may later
attach an estimated or measured probability distribution without changing the set
of legal traces.

\section{Scenario 3: exchange the code and use the token}

\begin{lstlisting}
def exchangeAndFetch : Scenario :=
  scenario "exchange code and fetch resource" do
    use code     <- imported AuthCode "code"
    use resource <- imported Resource "resource"
    let token    <- declare AccessToken "accessToken"
    let profile  <- declare Profile "profile"

    from_ .app
    to_ .authorizationServer
    emit "token requested" [field "code" code]

    from_ .authorizationServer
    to_ .app
    emit "access token issued" [field "accessToken" token]

    from_ .app
    to_ .resourceServer
    emit "protected resource requested"
      [field "accessToken" token, field "resource" resource]

    from_ .resourceServer
    to_ .app
    emit "protected resource returned"
      [field "resource" resource, field "profile" profile]

    from_ .app
    to_ .browser
    emit "application response" [field "profile" profile]
\end{lstlisting}

The token itself is visible only at the endpoints where the abstract protocol says
it is transmitted. The grammar makes no claim about the authorization server's
database schema or the application's session representation.

\section{Reading the scenarios together}

\begin{lstlisting}
def oauthAuthorization (resource : Resource) (password : Password) : Scenario :=
  scenario "OAuth authorization code flow" do
    include (beginAuthorization resource)
    include (authenticateAtAuthorizationServer password)
    include authorizationDecision
    after authorizationGranted include exchangeAndFetch

def oauthSystem (resource : Resource) (password : Password) : SystemSpec :=
  generatedBy (oauthAuthorization resource password)
  |> requireAll oauthRequirements
\end{lstlisting}

The \texttt{after} line is a guarded nonterminal reference. The denial branch ends
without exchanging a code. All scenario sources are read into one grammar, but
their names and boundaries remain available to render and test independently.

\section{Interaction diagrams by scenario}

\begin{figure}[p]
\centering
\resizebox{\textwidth}{!}{\OAuthInteractionDiagram{black}{white}{black!4}}
\caption{OAuth interaction view. The frames bracket the two alternatives after
the common authorization prefix; only the granted branch continues to token
exchange and resource access.}
\label{fig:oauth-interaction}
\end{figure}

\section{Requirements as observable constraints}

\begin{lstlisting}
def oauthRequirements : List Requirement :=
  [ require "callback state matches request"
      (on "authorization callback" (previouslySame "state"
        "authorization requested"))

  , forbid "code exchange after denial"
      (previously "authorization denied" && event "token requested")

  , require "issued token came from a code"
      (on "access token issued" (previouslySame "code" "token requested"))

  , ctl "every decision reaches the application"
      (AG (decisionAtAuthorizationServer --> AF callbackAtApplication))

  , ctl "a protected result is never returned without a token request"
      (AG (resourceReturned --> previously tokenRequested))
  ]
\end{lstlisting}

\subsection{Password secrecy}

\index{password secrecy}
Observable events are not the same as universally readable plaintext. Every actor
may observe that credential bytes were sent; only actors within the trusted
user-agent boundary and authorization service may derive the password value. The
theorem must name both the observer and its assumptions.

\begin{lstlisting}
theorem oauth_password_secret
    (trace : Trace WireEvent)
    (accepted : (oauthSystem resource password).Accepts trace)
    (initial : InitiallyKnown password [.browser, .authorizationServer])
    (honest : Uncompromised [.browser, .authorizationServer])
    (protection : PerfectProtection protectFor) :
    not (Derives (knowledgeOf .networkAttacker trace) (secret password)) := by
  exact protectedFieldNotDerivable accepted initial honest protection
\end{lstlisting}

This statement deliberately permits the trusted browser boundary and authorization
service to know the password. A theorem claiming secrecy from the browser would be
false for a password typed into that browser. If the desired boundary is narrower,
the actor model must represent a trusted credential component separately.

The theorem is also conditional. A compromised endpoint, logging side channel, or
failed concrete protection refinement invalidates an assumption rather than being
hidden by the grammar.

\subsection{Byte-level refinement without choosing a transport}

Every vague OAuth message can be assigned a canonical byte term before HTTP or
protobuf is chosen:

\begin{lstlisting}
refineMessage "authorization requested" where
  term := tuple
    [ tagBytes 0x01
    , namedBytes "state" (encode state)
    , namedBytes "redirectUri" (encode callback)
    ]
  bytes := encodeCanonical term

refineMessage "token requested" where
  term := tuple [tagBytes 0x02, namedBytes "code" (encode code)]
  bytes := encodeCanonical term
\end{lstlisting}

This fixes framing, tags, and value relationships at an abstract byte-algebra level.
A subsequent HTTP mapping might put those bytes into form fields; a protobuf mapping
might use numbered fields. Either concrete alphabet must prove that decoding and
re-encoding preserve the abstract terminal and its parameter bindings.

The use of \texttt{previously} is deliberately a trace predicate in this first
version. It can later receive dedicated past-time syntax without changing the event
model.

\section{Probabilities and charts}

An operational model may refine the decision and exchange outcomes:

\begin{lstlisting}
def measuredAuthorizationDecision : Scenario :=
  choose "authorization decision"
    [ weighted (87 / 100) authorizationGranted
    , weighted (13 / 100) authorizationDenied ]

def oauthCharts : List ChartSpec :=
  [ pie "authorization outcomes" do
      source events; groupBy terminalOutcome; value count
  , line "end-to-end authorization latency" do
      source (correlate events by "state")
      x tstamp; y (elapsed "authorization requested" "application response")
  ]
\end{lstlisting}

The percentages are model inputs or measurements, not truths inferred from the
grammar. The grammar continues to define which behaviours are possible; the
probabilistic refinement says how mass is distributed over them.

\section{Residual state machine}

\begin{figure}[htbp]
\centering
\resizebox{\textwidth}{!}{%
\begin{tikzpicture}[->,>=Latex,node distance=2.35cm,on grid,auto,
  every state/.style={draw,minimum size=.9cm,font=\small}]
  \node[state,initial] (start) {start};
  \node[state] (requested) [right=of start] {requested};
  \node[state] (decision) [right=of requested] {decision};
  \node[state] (callback) [above right=1.25cm and 2.5cm of decision] {callback};
  \node[state,accepting] (denied) [below right=1.25cm and 2.5cm of decision] {denied};
  \node[state] (token) [right=of callback] {token};
  \node[state] (fetch) [right=of token] {fetch};
  \node[state,accepting] (done) [right=of fetch] {done};
  \path
    (start) edge node {request} (requested)
    (requested) edge node {authorize} (decision)
    (decision) edge node {code} (callback)
    (decision) edge node[below] {deny} (denied)
    (callback) edge node {exchange} (token)
    (token) edge node {issued} (fetch)
    (fetch) edge node {resource} (done);
\end{tikzpicture}}
\caption{OAuth state machine derived from legal grammar continuations.}
\label{fig:oauth-state-machine}
\end{figure}

This automaton is intentionally behavioural. It does not contain browser cookies,
database rows, server threads, or cryptographic library state. Those details may be
essential to an implementation proof, but they are refinements of this contract,
not the contract itself.
